
\section{\trFloerHomology\trInvariance{}证明的收尾：命题~11.2.9~的证明}
\label{sec:conclusion-invariance}

设 $\Gamma'=(H',J')$ 与 $\Gamma''=(H'',J'')$ 是分别把 $(H^{a},J^{a})$ 连接到
$(H^{b},J^{b})$、把 $(H^{b},J^{b})$ 连接到 $(H^{c},J^{c})$ 的两条\trRegular\trHomotopy。正如本章前面各节中那样，这两条\trHomotopy{}被取成在 $|s|\ge R$ 时是\trStationary{}的，其中 $R>0$ 是一个尚未取定的常数（在证明的进程中它可以增大）。

$\Gamma'$ 与 $\Gamma''$ 的拼接借助公式
\[
H_{\rho}(s,t,p)=
\left\{\begin{array}{ll}
H'(s+\rho,t,p) & \text{若 } s\le 0\\
H''(s-\rho,t,p) & \text{若 } s\ge 0
\end{array}\right.
\]
\[
J_{\rho}(s,p)=
\left\{\begin{array}{ll}
J'(s+\rho,p) & \text{若 } s\le 0\\
J''(s-\rho,p) & \text{若 } s\ge 0\text{。}
\end{array}\right.
\]
定义了一条\trHomotopy{} $\Gamma_{\rho}=(H_{\rho},J_{\rho})$。

这条\trHomotopy{} $\Gamma_{\rho}$ 在 $\rho$ 充分大（$\rho>R$）时是有定义的，它把
$(H^{a},J^{a})$ 连接到 $(H^{c},J^{c})$，并且在 $|s|\ge R+\rho$ 时是\trStationary{}的。注意，\trApriori{}地看，它可能不是\trRegular{}的。我们将证明下面这个结果。

\begin{lemma}\label{thm:gamma-rho-regular-large}
在对 $\Gamma'$ 与 $\Gamma''$ 作一个小的\trPerturbation{}（在\trRegular\trHomotopy{}的范围内）之后，我们可以假定：存在任意大的 $\rho$，使得 $\Gamma_{\rho}$ 是\trRegular{}的。
\end{lemma}

于是存在一个趋于无穷的序列 $\rho_{n}$，使得 $\Gamma_{\rho_{n}}$ 对每个 $n$ 都是\trRegular{}的。

对于这些 $\rho$ 值，我们就可以定义\cref{sec:morphism-phi-gamma}中\trComplex{}的\trMorphism
\[
\Phi^{\Gamma_{\rho}}\colon CF_{\star}\bigl(H^{a},J^{a}\bigr)
\longrightarrow CF_{\star}\bigl(H^{c},J^{c}\bigr)\text{。}
\]

注意，由\cref{thm:homology-independent-gamma}，\cref{thm:gamma-rho-regular-large}中对 $\Gamma'$、$\Gamma''$ 的那些\trPerturbation{}在\trHomology{}层面上分别定义出与 $\Phi^{\Gamma'}$、$\Phi^{\Gamma''}$ 相同的\trMorphism。因此，对这些\trHomotopy{}我们将保留同样的记号 $\Gamma'$、$\Gamma''$。

\cref{thm:composition-regular-homotopy}显然是下面这个结果的推论。

\begin{proposition}\label{thm:exists-rho-coincide}
存在一个 $\rho$，使得 $\Gamma_{\rho}$ 是\trRegular{}的，并且\trComplex{}的诸\trMorphism{} $\Phi^{\Gamma''}\circ\Phi^{\Gamma'}$\footnote{译注：原书此处两个因子均带上标 $\Gamma''$（印作 $\Phi^{\Gamma''}\circ\Phi^{\Gamma''}$）；按 $\Phi^{\Gamma'}$ 与 $\Phi^{\Gamma''}$ 各自的定义域与陪域（并参见下文的证明），复合应作 $\Phi^{\Gamma''}\circ\Phi^{\Gamma'}$，疑系排印之误，此处径改。——译者注}与 $\Phi^{\Gamma_{\rho}}$ 相重合。
\end{proposition}

\begin{proof}[\cref{thm:gamma-rho-regular-large}的证明]
考虑\cref{subsec:transversality-gamma}中定义的空间 $\cC^{\infty}_{\varepsilon}$。设 $\cC^{\infty}_{\varepsilon,0}$ 是诸函数
\[
h\colon\RR\times S^{1}\times W\longrightarrow\RR
\]
中满足 $h(s,t,p)=0$（当 $|s|<\delta$，其中 $\delta$ 是任意一个正数）者所成的子空间。我们取定一个 $\rho_{1}$，它充分大，使得 $\Gamma_{\rho_{1}}$ 有定义。给定一个函数 $h\in\cC^{\infty}_{\varepsilon,0}$，我们让它与按公式
\[
h'(s,t,p)=
\left\{\begin{array}{ll}
h(s-\rho_{1},t,p) & \text{若 } s\le\rho_{1}\\
0 & \text{若 } s\ge\rho_{1}
\end{array}\right.
\]
\[
h''(s,t,p)=
\left\{\begin{array}{ll}
0 & \text{若 } s\le-\rho_{1}\\
h(s+\rho_{1},t,p) & \text{若 } s\ge-\rho_{1}
\end{array}\right.
\]
所定义的两个函数 $h'$、$h''\in\cC^{\infty}_{\varepsilon}$ 相对应，使得由 $(H'+h',J')$ 与 $(H''+h'',J'')$ 所定义的\trHomotopy{} $\Gamma_{\rho_{1}}$ 恰好就是 $(H_{\rho_{1}}+h,J_{\rho_{1}})$。从而，要使 $\Gamma_{\rho_{1}}$ 成为\trRegular{}的，只需对（更小的）\trPerturbation{}空间 $\cC^{\infty}_{\varepsilon,0}$ 重复\cref{subsec:transversality-gamma}中给出的\trTransversality{}证明即可。唯一的小改动出现在与\cref{thm:z-identically-zero}类似的那条引理的证明中。若\trVectorField{} $Z$ 非零，则存在一点 $(s_{0},t_{0})\in\bigl(\RR-\{0\}\bigr)\times S^{1}$，使得 $Z(s_{0},t_{0})\ne0$。证明的其余部分相同。

随后我们可以对一个 $\rho_{2}>\rho_{1}$ 重复这一证明，此时要选取充分小的\trPerturbation{} $h'$ 与 $h''$，使得 $\Gamma_{\rho_{1}}$ 保持\trRegular{}；于是 $\Gamma_{\rho_{1}}$ 与 $\Gamma_{\rho_{2}}$ 就都将是\trRegular{}的。以类似的方式继续下去，我们就构造出一个趋于无穷的序列 $\rho_{n}$，使得 $\Gamma_{\rho_{n}}$ 对每个 $n$ 都是\trRegular{}的。
\end{proof}

\begin{proof}[\cref{thm:exists-rho-coincide}的证明]
设 $x\in\operatorname{Crit}(\scr{A}_{H^{a}})$。采用\cref{subsec:compactness-gamma}中的记号，我们有
\[
\begin{aligned}
\Phi^{\Gamma''}\circ\Phi^{\Gamma'}(x)
&=\Phi^{\Gamma''}\Biggl(\sum_{\substack{y\in\operatorname{Crit}(\scr{A}_{H^{b}})\\ \mu(y)=\mu(x)}}
n^{\Gamma'}(x,y)\cdot y\Biggr)\\
&=\sum_{\substack{z\in\operatorname{Crit}(\scr{A}_{H^{c}})\\ \mu(z)=\mu(x)}}
\ \sum_{\substack{y\in\operatorname{Crit}(\scr{A}_{H^{b}})\\ \mu(y)=\mu(x)}}
n^{\Gamma'}(x,y)\,n^{\Gamma''}(y,z)\cdot z\text{。}
\end{aligned}
\]
如同\cref{sec:morphism-phi-gamma}中那样，$n^{\Gamma'}(x,y)$ 与 $n^{\Gamma''}(y,z)$ 分别是
$\scr{M}^{\Gamma'}(x,y)$ 与 $\scr{M}^{\Gamma''}(y,z)$ 中元素个数（模 $2$）。显然，\cref{thm:exists-rho-coincide}是下面这个结果的推论。
\end{proof}

\begin{proposition}\label{thm:chi-bijection}
设 $x\in\operatorname{Crit}(\scr{A}_{H^{a}})$ 与 $z\in\operatorname{Crit}(\scr{A}_{H^{c}})$ 满足 $\mu(x)=\mu(z)$。对每个充分大的 $\rho$，诸集合
\[
\bigcup_{\substack{y\in\operatorname{Crit}(\scr{A}_{H^{b}})\\ \mu(y)=\mu(x)}}
\scr{M}^{\Gamma'}(x,y)\times\scr{M}^{\Gamma''}(y,z)
\qquad\text{与}\qquad
\scr{M}^{\Gamma_{\rho}}(x,z)
\]
之间存在一个\trBijective。
\end{proposition}

我们如今来证明这条命题。以下是证明的思路。它归结为（这并不出人意料）：
\begin{enumerate}[label=(\arabic*)]
\item 对充分大的 $\rho$，定义一个映射
\[
\chi_{\rho}\colon
\bigcup_{\substack{y\in\operatorname{Crit}(\scr{A}_{H^{b}})\\ \mu(y)=\mu(x)}}
\scr{M}^{\Gamma'}(x,y)\times\scr{M}^{\Gamma''}(y,z)
\longrightarrow\scr{M}^{\Gamma_{\rho}}(x,z)
\]
\item 证明它对充分大的 $\rho$ 是\trInjective{}的。
\item 证明它对充分大的 $\rho$ 是\trSurjective{}的。
\end{enumerate}

为了对偶对 $(u,v)\in\scr{M}^{\Gamma'}(x,y)\times\scr{M}^{\Gamma''}(y,z)$ 定义 $\chi_{\rho}(u,v)$，我们如同\cref{sec:construction-psi-rho}与\cref{subsec:the-gluing}中那样使用一个\trGluing{}过程：从\trPregluing{} $w_{\rho}=u\,\#_{\rho}\,v$ 出发，我们借助\trNewton--\trPicard{}方法构造一个解 $\varphi(\rho)\in\scr{M}^{\Gamma_{\rho}}$。

至于用这一方法得到的其他那些\trFloer{}解，我们将发现：当 $\rho$ 趋于无穷时，$\varphi_{\rho}$ 趋于 $(u,v)$（在朝\trBrokenTrajectory{}\trConvergence{}的意义下；参见\cref{subsec:compactness-gamma}）。

要证明 $\chi$ 的\trInjective{}性\footnote{译注：原书此处印作“$\varphi$ 的\trInjective{}性”；按此处所要证明的对象与本段下文均为映射 $\chi_{\rho}$（参见上文证明纲要之 (2) 与小节 11.5.b 的标题），而 $\varphi(\rho)=\chi_{\rho}(u,v)$ 只是一个元素，疑应为 $\chi$。此处径改。——译者注}，我们使用反证法。若 $\chi_{\rho}$ 从 $\rho$ 的某个值起不是\trInjective{}的，那么就会存在一个趋于无穷的序列 $\rho_{n}$ 以及偶对 $(u_{n},v_{n})\ne(u'_{n},v'_{n})$，使得 $\chi_{\rho_{n}}(u_{n},v_{n})=\chi_{\rho_{n}}(u'_{n},v'_{n})$。由于诸序列 $(u_{n},v_{n})$ 在一个有限集中取值，所以在必要时抽取子序列之后，我们可以——而且我们也确实如此——假定它们是常值的，比如说分别等于 $(u,v)$ 与 $(u',v')$。在等式
\[
\chi_{\rho_{n}}(u,v)=\chi_{\rho_{n}}(u',v')
\]
中令 $n$ 趋于无穷，我们便得到 $(u,v)=(u',v')$，这是一个矛盾。

至于 $\chi_{\rho}$ 的\trSurjective{}性，推理是类似的。若对一个趋于无穷的序列 $\rho_{n}$，存在一个不在 $\chi_{\rho_{n}}$ 的像中的 $\ell_{n}$，那么我们先证明：$\lim\ell_{n}=(u,v)$，其中 $(u,v)$ 是 $\scr{M}^{\Gamma'}(x,y)\times\scr{M}^{\Gamma''}(y,z)$ 的某个偶对（$y\in\operatorname{Crit}\scr{A}_{H^{b}}$，$\mu(y)=\mu(x)$）。接着我们应用\trGluing{}的一条唯一性性质（类似于\cref{sec:uniqueness-gluing}中的那一条）来证明：对充分大的 $n$，$\ell_{n}\in\operatorname{Im}\chi_{n}$，从而导出一个矛盾。

以下是这一证明的细节。

\subsection{$\chi$ 的定义}
\label{subsec:definition-chi}

设 $y\in\operatorname{Crit}\scr{A}_{H^{b}}$ 满足 $\mu(y)=\mu(x)$。考虑 $\scr{M}^{\Gamma'}(x,y)\times\scr{M}^{\Gamma''}(y,z)$ 中的一个偶对 $(u,v)$。$u$ 与 $v$ 的\trPregluing{} $w_{\rho}$ 由与\cref{sec:pregluing}中相同的公式来定义。由于 $w_{\rho}(s,t)=u(s+\rho,t)$（若 $s\le-1$）而 $w_{\rho}(s,t)=v(s-\rho,t)$（若 $s\ge1$），$\Gamma_{\rho}=(H_{\rho},J_{\rho})$ 的定义便蕴含
\[
\Bigl(\frac{\partial w_{\rho}}{\partial s}
+J_{\rho}\frac{\partial w_{\rho}}{\partial t}
+\operatorname{grad}_{w_{\rho}}H_{\rho}\Bigr)(s,t)=0
\qquad\text{对 } |s|\ge 1\text{。}
\]
当 $\rho$ 趋于 $+\infty$ 时，这个表达式对于 $\cC^{\infty}_{\mathrm{loc}}$ 范数（从而也对 $L^{p}$ 范数）趋于 $0$。的确，
\[
\lim_{\rho\to+\infty}w_{\rho}(s,t)=y(t)
\qquad\text{与}\qquad
\lim_{\rho\to+\infty}\bigl(H_{\rho},J_{\rho}\bigr)=\bigl(H^{b},J^{b}\bigr)
\]
（两种情形都按 $\cC^{\infty}_{\mathrm{loc}}$ 范数理解）。

解 $\varphi(\rho)\in\scr{M}^{\Gamma_{\rho}}(x,z)$ 将从一个\trApproximateSolution{} $w_{\rho}$ 出发、借助一个形如下式的公式来得到：
\[
\varphi(\rho)=\exp_{w_{\rho}}\bigl(\gamma_{\rho}\bigr)\text{。}
\]
因此它将是
\[
Y\overset{\scr{F}_{\rho}}{\longmapsto}
\Bigl(\frac{\partial}{\partial s}+J_{\rho}\frac{\partial}{\partial t}
+\operatorname{grad}H_{\rho}\Bigr)\bigl(\exp_{w_{\rho}}Y\bigr)
\]
的一个零点。

利用恰如\cref{sec:construction-psi-rho}中那样构造的诸\trFrame{} $(Z_{i}^{\rho})$（位于 $\scr{F}_{\rho}$ 的源与靶处），我们可以——而且我们也确实如此——把 $\scr{F}_{\rho}$ 视为一个映射
\[
\scr{F}_{\rho}\colon W^{1,p}\bigl(\RR\times S^{1};\RR^{2n}\bigr)
\longrightarrow L^{p}\bigl(\RR\times S^{1};\RR^{2n}\bigr)\text{。}
\]
\trLinearization{}算子 $L_{\rho}=(d\scr{F}_{\rho})_{0}$ 具有形状
\[
Y\longmapsto\frac{\partial Y}{\partial s}
+J_{0}\frac{\partial Y}{\partial t}
+S_{\rho}\cdot Y\text{，}
\]
其中 $S_{\rho}$ 是一个从 $\RR\times S^{1}$ 到 $M_{2n}(\RR)$ 的映射。它是一个\trIndex{}为 $0$ 的\trFredholmOperator，因为 $\mu(x)=\mu(z)$。

\cref{thm:lrho-gamma-lower-bound}的一个类似物（它可以用完全相同的方式来证明）随后便蕴含：$L_{\rho}$ 是\trBijective{}的，并且具有一个逆映射
\[
G_{\rho}\colon L^{p}\bigl(\RR\times S^{1};\RR^{2n}\bigr)
\longrightarrow W^{1,p}\bigl(\RR\times S^{1};\RR^{2n}\bigr)
\]
其范数被一个不依赖于 $\rho$ 的常数所控制。我们可以应用\trNewton--\trPicard{}方法所需的那些条件（\cref{thm:abstract-newton-picard}）都是满足的：与\cref{thm:dn-op-estimate}中那条估计相类似、使我们得以验证\cref{thm:abstract-newton-picard}之条件 (2) 的那条估计，也可以按完全相同的方式来证明。

于是我们得到一个 $\gamma_{\rho}\in W^{1,p}\bigl(\RR\times S^{1};\RR^{2n}\bigr)$，使得
\[
\varphi(\rho)=\exp_{w_{\rho}}\gamma_{\rho}\in\scr{M}^{\Gamma_{\rho}}(x,z)\text{。}
\]
从而我们可以令 $\chi_{\rho}(u,v)=\varphi(\rho)$。如同\cref{thm:gamma-rho-tends-zero}中那样，我们得到
\[
\lim_{\rho\to+\infty}\bigl\|\gamma_{\rho}\bigr\|_{W^{1,p}}=0\text{。}
\]
恰如\cref{sec:construction-psi-rho}中那样，这就蕴含了下面的结果。

\begin{proposition}\label{thm:phirho-convergence}
元素 $\varphi(\rho)$ 满足
\[
\lim_{\rho\to+\infty}\varphi(\rho)(s-\rho,t)=u(s,t)
\qquad\text{与}\qquad
\lim_{\rho\to+\infty}\varphi(\rho)(s+\rho,t)=v(s,t)
\]
于 $\cC^{\infty}_{\mathrm{loc}}$ 中。\noqed
\end{proposition}

\subsection{$\chi$ 的\trInjective{}性}
\label{subsec:injectivity-chi}

我们假定 $\chi$ 不是\trInjective{}的，也就是说，存在任意大的 $\rho$，使得 $\chi_{\rho}$ 不是\trInjective{}的。考虑三个序列 $\rho_{n}$、$(u_{n},v_{n})$ 与 $(u'_{n},v'_{n})$，满足
\[
(u_{n},v_{n})\ \text{与}\ (u'_{n},v'_{n})\in
\bigcup_{\substack{y\in\operatorname{Crit}(\scr{A}_{H^{b}})\\ \mu(y)=\mu(x)}}
\scr{M}^{\Gamma'}(x,y)\times\scr{M}^{\Gamma''}(y,z)
\]
以及
\[
\lim_{n\to+\infty}\rho_{n}=+\infty\text{，}\qquad
(u_{n},v_{n})\ne(u'_{n},v'_{n})\text{，}\qquad
\chi_{\rho_{n}}(u_{n},v_{n})=\chi_{\rho_{n}}(u'_{n},v'_{n})\text{。}
\]

诸序列 $(u_{n},v_{n})$ 与 $(u'_{n},v'_{n})$ 在一个有限集中取值。因此，在必要时抽取子序列之后，我们可以——而且我们也确实如此——假定它们是常值的，比如说分别等于 $(u,v)$ 与 $(u',v')$，满足
\[
(u,v)\ne(u',v')
\qquad\text{且}\qquad
\chi_{\rho_{n}}(u,v)=\chi_{\rho_{n}}(u',v')\text{。}
\]
由\cref{thm:phirho-convergence}，我们如今有
\[
\begin{aligned}
u(s,t)&=\lim_{n\to+\infty}\bigl(\chi_{\rho_{n}}(u,v)\bigr)(s-\rho_{n},t)\\
&=\lim_{n\to+\infty}\bigl(\chi_{\rho_{n}}(u',v')\bigr)(s-\rho_{n},t)\\
&=u'(s,t)
\end{aligned}
\]
并且同样地有 $v(s,t)=v'(s,t)$。$\chi_{\rho}$ 的\trInjective{}性随之得证。

\subsection{$\chi$ 的\trSurjective{}性}
\label{subsec:surjectivity-chi}

我们假定（再一次为了导出矛盾）：存在任意大的 $\rho$，使得 $\chi_{\rho}$ 不是\trSurjective{}的。考虑一个趋于无穷的序列 $(\rho_{n})$ 以及一列元素 $\ell_{n}\in\scr{M}^{\Gamma_{\rho_{n}}}(x,z)$，其中 $\ell_{n}\notin\operatorname{Im}(\chi_{\rho_{n}})$（对每个 $n$）。我们将证明下面这条命题。

\begin{proposition}\label{thm:broken-limit-extract}
存在一点 $y\in\operatorname{Crit}\scr{A}_{H^{b}}$，它满足 $\mu(y)=\mu(x)$，以及一个偶对
\[
(u,v)\in\scr{M}^{\Gamma'}(x,y)\times\scr{M}^{\Gamma''}(y,z)
\]
使得，在必要时抽取一个子序列之后，
\[
\lim_{n\to+\infty}\ell_{n}(s-\rho_{n},t)=u(s,t)
\qquad\text{与}\qquad
\lim_{n\to+\infty}\ell_{n}(s+\rho_{n},t)=v(s,t)
\]
按 $\cC^{\infty}_{\mathrm{loc}}$ \trTopology{}成立。
\end{proposition}

为了证明，我们将需要下面这个结果。

\begin{proposition}\label{thm:subsequence-floer-solution}
\begin{enumerate}[label=(\arabic*)]
\item 设 $(s_{n})$ 是一个实数序列。那么序列 $\ell_{n}(s+s_{n},t)$ 承认一个按 $\cC^{\infty}_{\mathrm{loc}}$ \trTopology{}\trConvergence{}的子序列。
\item 设 $\ell$ 是这样一个子序列的极限。那么存在一个常数 $s_{\star}$，使得 $\ell(s+s_{\star},t)$ 是与诸偶对 $(H^{a},J^{a})$、$(H',J')$、$(H^{b},J^{b})$、$(H'',J'')$、$(H^{c},J^{c})$ 之一相对应的\trFloerEquation{}的一个解。
\end{enumerate}
\end{proposition}

\begin{proof}[\cref{thm:subsequence-floer-solution}的证明]
我们先来证明点 (2)，其中承认点 (1)。映射 $(s,t)\longmapsto\ell_{n}(s+s_{n},t)$ 是与偶对
\[
\bigl(H_{n},J_{n}\bigr)=\bigl(H_{\rho_{n}}(s+s_{n},t,p),J_{\rho_{n}}(s+s_{n},p)\bigr)
\]
相对应的\trFloerEquation{}的一个解。回想一下，按定义，$H_{n}$ 与 $J_{n}$ 由
\[
H_{n}(s,t,p)=
\left\{\begin{array}{ll}
H'(s+\rho_{n}+s_{n},t,p) & \text{若 } s\le-s_{n}\\
H''(s-\rho_{n}+s_{n},t,p) & \text{若 } s\ge-s_{n}
\end{array}\right.
\]
\[
J_{n}(s,p)=
\left\{\begin{array}{ll}
J'(s+\rho_{n}+s_{n},p) & \text{若 } s\le-s_{n}\\
J''(s-\rho_{n}+s_{n},p) & \text{若 } s\ge-s_{n}\text{。}
\end{array}\right.
\]
给出。让我们按照序列 $(s_{n})$ 的性态来研究 $(H_{n},J_{n})$ 的各个可能的极限（按 $\cC^{\infty}_{\mathrm{loc}}$ \trTopology）。

在 $s_{n}=-\rho_{n}$ 这一特殊情形，$(H_{n},J_{n})$ 显然趋于 $(H',J')$。同样地，当 $s_{n}=\rho_{n}$ 时，$(H_{n},J_{n})$ 趋于 $(H'',J'')$。于是，极限 $\ell$ 在第一种情形将是与 $\Gamma'=(H',J')$ 相对应、在第二种情形将与 $\Gamma''=(H'',J'')$ 相对应的\trFloerEquation{}的一个解。

由于点 (2) 的陈述允许我们给 $(s_{n})$ 添加一个\trConvergence{}的序列，所以我们只需研究诸序列 $(s_{n}\pm\rho_{n})$ 皆发散的情形：
\begin{itemize}
\item 若 $\lim(s_{n}+\rho_{n})=-\infty$，也就是说，若 $s_{n}\ll-\rho_{n}$，那么我们容易得到
\[
\lim_{n\to+\infty}\bigl(H_{n},J_{n}\bigr)=\bigl(H^{a},J^{a}\bigr)\text{。}
\]
\item 若 $\lim_{n\to+\infty}(s_{n}+\rho_{n})=+\infty$ 且 $\lim_{n\to+\infty}(s_{n}-\rho_{n})=-\infty$，换句话说，若 $-\rho_{n}\ll s_{n}\ll\rho_{n}$，那么通过把 $(s_{n})$ 趋于 $-\infty$、\trConvergence{}与趋于 $+\infty$ 这几种情形分开来分析，我们发现
\[
\lim_{n\to+\infty}\bigl(H_{n},J_{n}\bigr)=\bigl(H^{b},J^{b}\bigr)\text{。}
\]
\item 若 $\lim(s_{n}-\rho_{n})=+\infty$，也就是说，若 $s_{n}\gg\rho_{n}$，那么我们有
\[
\lim_{n\to+\infty}\bigl(H_{n},J_{n}\bigr)=\bigl(H^{c},J^{c}\bigr)\text{。}
\]
\end{itemize}
这样我们就证明了点 (2)；现在让我们转向点 (1)。我们先照抄\cref{thm:uniform-gradient-bound}（\trCompact{}性定理的证明中对\trGradient{}的一个界）的证明，来证明存在一个常数 $A>0$，使得
\[
\forall\rho,\ \forall\ell\in\scr{M}^{\Gamma_{\rho}}(x,z),\
\forall(s,t)\in\RR\times S^{1},\quad
\bigl\|\operatorname{grad}_{(s,t)}\ell\bigr\|\le A
\]
（这里 $\|\cdot\|$ 表示 $TW$ 上的任意一个范数——由于 $W$ 是\trCompact{}的，这些范数都是等价的——比如说从 $\RR^{m}$ 承来的那个范数）。

若这样的常数不存在，那么如同\cref{thm:uniform-gradient-bound}的证明中那样，我们就会得到一条 $J$-\trHolomorphic{}曲线 $v$，它将如同\cref{thm:symplectic-area-nonzero}中那样导出一个矛盾。此处所说的结构 $J$ 是序列 $J_{\rho_{n}}$ 的（按 $\cC^{\infty}_{\mathrm{loc}}$ \trTopology{}的）极限。因此它要么等于某个 $\rho_{\star}>0$ 的 $J_{\rho_{\star}}$（此时它依赖于 $s$），要么等于 $J^{b}$。于是，上面的不等式对每个 $n$ 的 $\ell_{n}$ 都成立，并且对序列 $\bigl(\ell_{n}(s+s_{n},t)\bigr)$ 的诸元素也成立。后者构成一个\trEquicontinuity{}的族，这就使得我们能够应用\trArzelaAscoliTheorem，从而找到一个按 $\cC^{0}_{\mathrm{loc}}$ \trTopology{}\trConvergence{}到极限 $\ell$ 的子序列。

点 (2) 随后表明：存在一个常数 $s_{\star}\in\RR$，使得 $\ell(s+s_{\star})$ 是（按弱意义）该命题中所陈述的诸\trFloerEquation{}之一的一个解。\trEllipticRegularity（\cref{thm:elliptic-regularity-lemma}）随后便蕴含 $\ell$ 是 $\cC^{\infty}$ 类的，并且 $\ell_{n}(s+s_{n},t)$ 朝 $\ell$ 的\trConvergence{}的确按 $\cC^{\infty}_{\mathrm{loc}}$ 意义成立。

至此\cref{thm:subsequence-floer-solution}的证明结束。
\end{proof}

为了结束\cref{thm:broken-limit-extract}的证明，我们将使用一个朝\trBrokenTrajectory{}\trConvergence{}的论证，它类似于我们在\cref{subsec:broken-trajectories}中使用过的那一个。更确切地说，一个类似于\cref{thm:broken-limit-decomposition}的证明给出了诸序列 $(s_{n}^{k})$（对 $0\le k\le m$）的存在性，使得（在必要时抽取一个子序列之后）对每个 $k$，序列 $\ell_{n}(s+s_{n}^{k},t)$ 都\trConvergence{}到一个极限 $\ell\in\cC^{\infty}(x_{k},x_{k+1})$，其中 $x_{0}=x$ 而 $x_{m+1}=z$。\cref{thm:subsequence-floer-solution}的点 (2) 表明，在必要时给诸序列 $(s_{n}^{k})_{n}$ 添加一些常数之后，诸极限 $\ell^{k}$ 都是与诸偶对 $(H^{a},J^{a})$、$(H',J')$、$(H^{b},J^{b})$、$(H'',J'')$、$(H^{c},J^{c})$ 之一相对应的\trFloerEquation{}的解。从而诸\trLoop{} $x_{k}$ 都是 $\scr{A}_{H^{a}}$、$\scr{A}_{H^{b}}$ 或 $\scr{A}_{H^{c}}$ 的\trCriticalPoint。

另一方面，若存在一个极限 $\ell_{k}$——它是与 $(H^{a},J^{a})$、$(H^{b},J^{b})$ 或 $(H^{c},J^{c})$ 相关联的\trFloerEquation{}的一个关于 $s$ 非常值的解——那么我们可以推出 $\mu(x_{k})-\mu(x_{k+1})\ge1$（如同\cref{thm:maslov-index-lower-bound}中那样），从而 $\mu(x)-\mu(z)\ge1$，这与我们所做的假设（$\mu(x)=\mu(z)$）相矛盾。

\begin{remark}\label{thm:constant-solutions-allowed}
注意，我们无法排除\trFloerEquation{}的那些关于 $s$ 为常值的解，因为我们无法排除诸\trHomotopy{} $\Gamma'$、$\Gamma''$ 或 $\Gamma_{\rho}$ 之一是平凡的情形。

也可能会碰巧诸偶对 $(H^{a},J^{a})$、$(H^{b},J^{b})$、$(H^{c},J^{c})$ 中的某些相等（而连接诸偶对的\trHomotopy{}却未必是平凡的）。因此我们需要考虑 $x$、$y$ 与 $z$ 不全相异的情形。
\end{remark}

让我们回到我们的证明。假定在诸极限 $\ell^{1},\dotsc,\ell^{m}$ 中有一个、比如说 $\ell^{k}$，它不是关于 $s$ 常值的。依照\cref{thm:subsequence-floer-solution}的点 (2) 的证明，我们可以——而且我们也确实如此——假定（在必要时给序列 $(s_{n}^{k})_{n}$ 添加一个常数之后）$\ell^{k}$ 是与 $\Gamma'$ 或 $\Gamma''$ 相关联的\trFloerEquation{}的一个解，并且 $s_{n}^{k}=-\rho_{n}$ 或 $s_{n}^{k}=\rho_{n}$。另一方面，\cref{thm:broken-limit-decomposition}的证明给出了具有如下性质的诸序列 $s_{n}^{k}$：
\[
\lim_{n\to+\infty}s_{n}^{k+1}-s_{n}^{k}=+\infty\text{。}
\]

因此，至多只有两个如上的极限 $\ell^{k}$，一个对应于 $s_{n}^{k}=-\rho_{n}$，另一个对应于 $s_{n}^{k}=\rho_{n}$：的确，这样的序列与 $-\rho_{n}$ 或与 $\rho_{n}$ 之差将是\trBounded{}的（由\cref{thm:subsequence-floer-solution}的点 (2)；回忆一下，不存在关于 $s$ 非常值、且是与 $(H^{a},J^{a})$、$(H^{b},J^{b})$ 或 $(H^{c},J^{c})$ 相关联的\trFloer{}解的 $\ell^{k}$）。

于是，序列 $\ell_{n}$ 所\trConvergence{}到的那条\trBrokenTrajectory{}由\trFloerEquation{}的一些关于 $s$ 为常值的解、至多一条与 $\Gamma'$ 相关联的\trTrajectory{}（对应于序列 $s_{n}=-\rho_{n}$）以及至多一条与 $\Gamma''$ 相关联的\trTrajectory{}（对应于 $s_{n}=\rho_{n}$）组成。再由\cref{thm:subsequence-floer-solution}，若有\footnote{译注：原书此式第一项 $\lim$ 的下标印作 $s\to+\infty$；按 $\ell_{n}$ 是一个序列而 $s$ 是自由变量（同一式中的第二项即作 $n\to+\infty$），疑应为 $n\to+\infty$。此处径改。——译者注}
\[
u=\lim_{n\to+\infty}\ell_{n}(s-\rho_{n},t)
\qquad\text{与}\qquad
v=\lim_{n\to+\infty}\ell_{n}(s+\rho_{n},t)\text{，}
\]
则
\[
u\in\scr{M}^{\Gamma'}(x',y')
\qquad\text{与}\qquad
v\in\scr{M}^{\Gamma''}(y'',z'')
\]
其中
\[
x\in\operatorname{Crit}\scr{A}_{H^{a}}\text{，}\qquad
y',y''\in\operatorname{Crit}\scr{A}_{H^{b}}\text{，}\qquad
z''\in\operatorname{Crit}\scr{A}_{H^{c}}\text{。}
\]
我们要证明的是：$x'=x$、$y'=y''$ 且 $z''=z$。

若在 $\ell_{n}$ 的极限的那条\trBrokenTrajectory{}中，有两条关于 $s$ 非常值的\trTrajectory{} $\ell'\in\cC^{\infty}(x,y)$ 与 $\ell''\in\cC^{\infty}(y,z)$，那么由上面的论述，$\ell'=u$ 且 $\ell''=v$，从而
\[
u\in\scr{M}^{\Gamma'}(x,y)
\qquad\text{与}\qquad
v\in\scr{M}^{\Gamma''}(y,z)\text{，}
\]
这就结束了证明。的确，不会有多于两条在 $s$ 方向上为常值的\trTrajectory\footnote{译注：原书此句印作“不会有多于两条在 $s$ 方向上为常值的\trTrajectory”；按上一段的论证（与 $\Gamma'$、$\Gamma''$ 相关联的非常值\trTrajectory{}各至多一条），疑应为“非常值”——若照字面取“常值”，则 $\ell'$ 之前、$\ell'$ 与 $\ell''$ 之间、$\ell''$ 之后均可有任意多条常值\trTrajectory，且“$\ell'$ 的到达点与 $\ell''$ 的出发点相重合”这一依据也无从给出常值者的上界。此处照原书排印，存疑备案。——译者注}，因为 $\ell'$ 的到达点与 $\ell''$ 的出发点相重合。

现在让我们假定：在 $\ell_{n}$ 的极限中只有一条关于 $s$ 非常值的\trFloer\trTrajectory{} $\ell\in\cC^{\infty}(x,z)$。这条\trTrajectory{} $\ell$ 于是等于 $u$ 或者 $v$。我们可以——而且我们也确实如此——不失一般性地假定 $\ell=u$。那么
\[
\lim_{n\to+\infty}\ell_{n}(s-\rho_{n},t)=u\in\scr{M}^{\Gamma'}(x,z)\text{。}
\]
利用关于\trBrokenTrajectory{}的\cref{thm:broken-limit-decomposition}的证明中的记号，存在一个 $s^{\star}\in\RR$，使得\trLoop{} $u(s,\cdot)\in B(z,\varepsilon)$（对 $s>s^{\star}$），从而对充分大的 $n$，\trLoop{} $\ell_{n}(s^{\star}-\rho_{n},\cdot)$ 也在球 $B(z,\varepsilon)$ 之中。

\cref{thm:broken-limit-decomposition}的证明随后蕴含：对每个 $s>s^{\star}-\rho_{n}$，都有 $\ell_{n}(s,\cdot)\in B(z,\varepsilon)$。否则，在极限中就会有另一条非常值的\trTrajectory{}（如本证明中那样构造）。这就蕴含：对 $v(s,t)=\lim_{n\to+\infty}\ell_{n}(s+\rho_{n},t)$，有 $v\in\scr{M}^{\Gamma''}(z,z)$（事实上，只要我们选取 $\varepsilon$ 充分小，这条\trTrajectory{}就是常值的且等于 $z$）。至此，在这种情形下也结束了\cref{thm:broken-limit-extract}的证明。

\begin{remark}\label{thm:limit-not-unique}
这里的困难来自这样一个事实：我们对\trBrokenTrajectory{}的\trConvergence{}所拥有的唯一定义（由\cref{thm:broken-limit-decomposition}给出）并不保证唯一的极限。\footnote{与\cref{thm:broken-limit-decomposition}所处理的情形相反。}因此我们必须求助于这条定理的证明，以构造这里所考虑的极限 $(\ell^{1},\dotsc,\ell^{m})$。
\end{remark}

最后，若 $\ell_{n}$ 的极限中的所有\trTrajectory{}都是常值的——在 $\cC^{\infty}(x,z)$ 中（特别地，此时我们有 $x=z$）——那么\cref{thm:broken-limit-decomposition}的证明表明：对每个 $s$ 都有 $\ell_{n}(s)\in B(x,\varepsilon)$。特别地，
\[
u(s,t)=\lim_{n\to+\infty}\ell_{n}(s-\rho_{n},t)\in\scr{M}^{\Gamma'}(x,z)
\]
而
\[
v(s,t)=\lim_{n\to+\infty}\ell_{n}(s+\rho_{n},t)\in\scr{M}^{\Gamma''}(x,z)\text{。}
\]

我们也可以证明 $u$ 是常值的且等于 $x$，$v$ 亦然，只要选取 $\varepsilon$ 充分小，使得 $\scr{M}^{\Gamma'}(x,x)$ 与 $\scr{M}^{\Gamma''}(x,x)$ 中所有关于 $s$ 非常值的\trTrajectory{}都穿出 $B(x,\varepsilon)$。

至此\cref{thm:broken-limit-extract}的证明结束。\noqed

我们如今有一个序列 $\ell_{n}\notin\operatorname{Im}(\chi_{\rho_{n}})$，使得
\[
\lim_{n\to+\infty}\ell_{n}(s-\rho_{n},t)=u(s,t)\in\scr{M}^{\Gamma'}(x,y)
\]
而
\[
\lim_{n\to+\infty}\ell_{n}(s+\rho_{n},t)=v(s,t)\in\scr{M}^{\Gamma''}(y,z)\text{。}
\]
为了得到所要的矛盾，我们来证明下面的结果。

\begin{proposition}\label{thm:elln-chi-rho}
对充分大的 $n$，$\ell_{n}=\chi_{\rho_{n}}(u,v)$。
\end{proposition}

这个证明沿着\trGluing{}的唯一性（\cref{thm:gluing-gamma}）之证明的主要线索进行。我们来证明\cref{thm:exists-rho-n-yn}的一个类似物。

\begin{proposition}\label{thm:elln-exp-yn}
对充分大的 $n$，存在一个 $Y_{n}\in w_{\rho_{n}}^{\star}TW$，使得
\[
\ell_{n}=\exp_{w_{\rho_{n}}}Y_{n}
\qquad\text{且}\qquad
\lim_{n\to+\infty}\bigl\|Y_{n}\bigr\|_{\infty}=0\text{。}
\]
\end{proposition}

在那之后，我们将证明\cref{thm:yn-w1p-null}的一个类似物。

\begin{proposition}\label{thm:yn-w1p-vanish}
\trVectorField{} $Y_{n}$ 属于 $W^{1,p}\bigl(w_{\rho_{n}}^{\star}TW\bigr)$。此外，$\lim_{n\to+\infty}\bigl\|Y_{n}\bigr\|_{W^{1,p}}=0$。
\end{proposition}

最后，我们将利用\trNewton--\trPicard{}方法所给出的唯一性（\cref{thm:abstract-newton-picard}）来结束证明。

\begin{proof}[\cref{thm:elln-exp-yn}的证明]
我们从一个引理开始。

\begin{lemma}\label{thm:elln-loc-uniform-y}
序列 $\bigl(\ell_{n}(s,t)\bigr)_{n}$ 按 $\cC^{\infty}_{\mathrm{loc}}$ \trTopology{}\trConvergence{}到 $y(t)$。
\end{lemma}

\begin{proof}
注意，对 $s\in[-\rho_{n}+R,\rho_{n}-R]$，我们有
\[
\bigl(H_{\rho_{n}},J_{\rho_{n}}\bigr)=\bigl(H^{b},J^{b}\bigr)\text{。}
\]
从而，$\ell_{n}$ 限制在 $[-\rho_{n}+R,\rho_{n}-R]\times S^{1}$ 上是与 $(H^{b},J^{b})$ 相关联的\trFloerEquation{}的一个解。特别地，函数
\[
s\longmapsto\scr{A}_{H^{b}}\bigl(\ell_{n}(s)\bigr)
\]
在这个区间上是递减的。

设 $a$ 表示 $\ell_{n}$ 的一个\trConvergence{}子序列的极限，其存在性由\cref{thm:subsequence-floer-solution}的断言 (1) 得出。这同一条命题的断言 (2) 的证明蕴含
\[
a\in\scr{M}^{(H^{b},J^{b})}\bigl(y',y''\bigr)
\qquad\text{对 } y',y''\in\operatorname{Crit}\scr{A}_{H^{b}}\text{。}
\]
我们取定一个 $s>R$ 与一个 $\sigma<-R$。对充分大的 $n$，我们有 $\sigma+\rho_{n}>s$。由此得出
\[
\begin{aligned}
\scr{A}_{H^{b}}\bigl(a(s)\bigr)
&=\lim_{n\to+\infty}\scr{A}_{H^{b}}\bigl(\ell_{n}(s)\bigr)\\
&\ge\lim_{n\to+\infty}\scr{A}_{H^{b}}\bigl(\ell_{n}(\sigma+\rho_{n})\bigr)\\
&=\scr{A}_{H^{b}}\bigl(v(\sigma)\bigr)\text{。}
\end{aligned}
\]
令 $s$ 趋于 $+\infty$、$\sigma$ 趋于 $-\infty$，我们发现
\[
\scr{A}_{H^{b}}\bigl(y''\bigr)\ge\scr{A}_{H^{b}}\bigl(y\bigr)\text{。}
\]
同样地，
\[
\begin{aligned}
\scr{A}_{H^{b}}\bigl(a(\sigma)\bigr)
&=\lim_{n\to+\infty}\scr{A}_{H^{b}}\bigl(\ell_{n}(\sigma)\bigr)\\
&\le\lim_{n\to+\infty}\scr{A}_{H^{b}}\bigl(\ell_{n}(s-\rho_{n})\bigr)\\
&=\scr{A}_{H^{b}}\bigl(u(s)\bigr)\text{。}
\end{aligned}
\]
我们再一次令 $s$ 趋于 $+\infty$、$\sigma$ 趋于 $-\infty$，便发现
\[
\scr{A}_{H^{b}}\bigl(y'\bigr)\le\scr{A}_{H^{b}}\bigl(y\bigr)\text{。}
\]
但是 $\scr{A}_{H^{b}}(y')\ge\scr{A}_{H^{b}}(y'')$，于是
\[
\scr{A}_{H^{b}}\bigl(y'\bigr)=\scr{A}_{H^{b}}\bigl(y''\bigr)=\scr{A}_{H^{b}}\bigl(y\bigr)\text{，}
\]
从而 $a$ 是关于 $s$ 常值的（更确切地说，$a(s)=y$），这是因为 $\scr{A}_{H^{b}}$ 的诸\trCriticalValue{}被假定是两两不同的。
\end{proof}

假定解 $u$ 与 $v$ 都不是关于 $s$ 常值的。我们在\trLoopSpace{} $\scr{L}W$ 中选取满足\cref{thm:exp-convex-ball}之性质的球 $B(x,\delta)$、$B(y,\delta)$ 与 $B(z,\delta)$。在必要时缩小 $\delta$，我们可以——而且我们也确实如此——还假定\trLoop{} $u(0)$ 与 $v(0)$ 不含于这些球中的任何一个。这些球除 $x$、$y$ 与 $z$ 之外，不含 $\operatorname{Crit}\scr{A}_{H^{a}}$、$\operatorname{Crit}\scr{A}_{H^{b}}$ 与 $\operatorname{Crit}\scr{A}_{H^{c}}$ 的任何其他\trLoop。

由刚才证明的那条引理，我们知道：对充分大的 $n$，$\ell_{n}(0)\in B(y,\delta)$。\footnote{译注：原书此句印作 $\ell(0)\in B(y,\delta)$（p.448）；按刚证引理 11.5.12（\cref{thm:elln-loc-uniform-y}，$\ell_{n}$ 按 $\cC^{\infty}_{\mathrm{loc}}$ \trTopology{}\trConvergence{}到 $y(t)$）可推知应为 $\ell_{n}(0)\in B(y,\delta)$，疑漏下标 $n$。此处径改。——译者注}如同\cref{thm:lift-exp-decay}的证明中那样，令
\[
\begin{aligned}
\sigma_{n}&=\sup\bigl\{s>0\ \bigm|\ \ell_{n}\bigl([-s,0]\bigr)\subset B(y,\delta)\bigr\}\\
\tau_{n}&=\sup\bigl\{s>0\ \bigm|\ \ell_{n}\bigl([0,s]\bigr)\subset B(y,\delta)\bigr\}\\
\sigma'_{n}&=\inf\bigl\{s>0\ \bigm|\ \ell_{n}\bigl(\,]{-\infty},-s]\bigr)\subset B(x,\delta)\bigr\}\\
\tau'_{n}&=\inf\bigl\{s>0\ \bigm|\ \ell_{n}\bigl([s,+\infty[\bigr)\subset B(z,\delta)\bigr\}\text{。}
\end{aligned}
\]
由于 $\lim\ell_{n}(-\rho_{n})=u(0)$ 且 $\lim\ell_{n}(\rho_{n})=v(0)$，我们有
\[
-\sigma'_{n}<-\rho_{n}<-\sigma_{n}
\qquad\text{与}\qquad
\tau_{n}<\rho_{n}<\tau'_{n}\text{。}
\]
这些序列的构造还蕴含\footnote{译注：原书此式第二项印作 $\ell_{n}(\sigma'_{n})\in\partial B(x,\delta)$；按 $\sigma'_{n}$ 的定义（$\ell_{n}$ 在 $]{-\infty},-s]$ 上落入 $B(x,\delta)$ 中的最后时刻是 $-\sigma'_{n}$），并且下面\cref{thm:four-differences-bounded}的证明也要用到 $\ell(0)=\lim_{n\to+\infty}\ell_{n}(-\sigma'_{n})\in\partial B(x,\delta)$，疑应为 $\ell_{n}(-\sigma'_{n})$。此处径改。——译者注}
\[
\ell_{n}(\sigma_{n}),\ \ell_{n}(\tau_{n})\in\partial B(y,\delta)\text{，}\quad
\ell_{n}(-\sigma'_{n})\in\partial B(x,\delta)\text{，}\quad
\ell_{n}(\tau'_{n})\in\partial B(z,\delta)\text{。}
\]
从而我们有下面的结果。

\begin{lemma}\label{thm:four-differences-bounded}
诸序列 $(\rho_{n}-\sigma_{n})$、$(\rho_{n}-\sigma'_{n})$、$(\rho_{n}-\tau_{n})$ 与 $(\rho_{n}-\tau'_{n})$ 都是\trBounded{}的。
\end{lemma}

\begin{proof}
假定 $(\sigma'_{n}-\rho_{n})$ 趋于 $+\infty$。那么由\cref{thm:subsequence-floer-solution}，序列 $\ell_{n}(s-\sigma'_{n})$ \trConvergence{}到 $\ell$——它是与 $(H^{a},J^{a})$ 相关联的\trFloerEquation{}的一个解（必要时参见\cref{thm:subsequence-floer-solution}的点 (2) 的证明）。此外，$\sigma'_{n}$ 的定义蕴含：对所有的负的 $s$ 都有 $\ell(s)\in B(x,\delta)$，从而随之得到
\[
\ell\in\scr{M}^{(H^{a},J^{a})}(x,y')
\qquad\text{对某个 } y'\in\operatorname{Crit}\scr{A}_{H^{a}}\text{。}
\]
由于 $\ell(0)\in\partial B(x,\delta)$，我们知道 $y'\ne x$。但这是不可能的，因为我们已经证明过（在证明\cref{thm:broken-limit-extract}时）：$\ell_{n}$ 所趋向的那条\trBrokenTrajectory{}不可能含有与 $(H^{a},J^{a})$ 相关联的\trFloerEquation{}的非常值\trTrajectory。

接下来，假定 $(\sigma_{n}-\rho_{n})$ 趋于 $-\infty$。仍由\cref{thm:subsequence-floer-solution}，序列 $\bigl(\ell_{n}(s-\sigma_{n})\bigr)$ \trConvergence{}到一个极限 $\ell$，它是与 $(H^{b},J^{b})$ 相关联的\trFloerEquation{}的一个解。我们有
\[
\ell(0)=\lim_{n\to+\infty}\ell_{n}(\sigma_{n})\in\partial B(y,\delta)\text{。}
\]
此外，$\sigma_{n}$ 趋于 $+\infty$（否则，若 $\sigma_{n}$ 是\trBounded{}的，那么由于 $\ell_{n}$ 按 $\cC^{\infty}_{\mathrm{loc}}$ \trConvergence{}到 $y$，$\ell_{n}(\sigma_{n})$ 就不可能含于 $\partial B(y,\delta)$ 之中）。

设 $s>0$。对充分大的 $n$，我们有 $s-\sigma_{n}\in[-\sigma_{n},0]$，于是 $\ell_{n}(s-\sigma_{n})\in B(y,\delta)$，从而 $\ell(s)\in B(y,\delta)$。由此得出：$\ell$ 是与 $(H^{b},J^{b})$ 相关联的\trFloerEquation{}的一个非常值解。它出现在 $\ell_{n}$ 的极限的那条\trBrokenTrajectory{}之中（这条\trTrajectory{}将以 $u,\ell$ 开始）。这是不可能的，理由同前。

我们以同样的方式进行，来证明关于 $\tau_{n}$ 与 $\tau'_{n}$ 的诸断言。
\end{proof}

设 $A>0$，使得前一陈述中的所有序列都在 $[-A,A]$ 中取值，并且使得
\[
\begin{aligned}
u(s)&\in B(x,\delta)\ \text{对 } s<-A\text{，}\qquad
u(s)\in B(y,\delta)\ \text{对 } s>A\text{，}\\
v(s)&\in B(y,\delta)\ \text{对 } s<-A
\qquad\text{且}\qquad
v(s)\in B(z,\delta)\ \text{对 } s>A\text{。}
\end{aligned}
\]
如同\cref{thm:lift-exp-decay}的证明中那样，我们把 $\RR$ 切成若干个区间，并对这些区间中的每一个定义 $Y_{n}(s)$：
\begin{enumerate}[label=(\arabic*)]
\item 若 $s<-\rho_{n}-A<-\sigma'_{n}$，那么 $\ell_{n}(s)\in B(x,\delta)$，
\[
w_{\rho_{n}}(s)=u(s+\rho_{n})\in B(x,\delta)
\]
并且我们使用\cref{thm:exp-convex-ball}。
\item 若 $-\rho_{n}-A<s<-\rho_{n}+A$，那么我们利用这一事实：$\ell_{n}(s-\rho_{n})$ 对 $s\in[-A,A]$ \trUniform{}地趋于 $u(s)$。
\item 若 $-\rho_{n}+A<s<\rho_{n}-A$，那么 $-\sigma_{n}<s<\tau_{n}$，从而 $\ell_{n}(s)\in B(y,\delta)$。再者，
\[
u(s+\rho_{n})
\qquad\text{与}\qquad
v(s-\rho_{n})\in B(y,\delta)\text{，}
\]
从而 $w_{\rho_{n}}(s)\in B(y,\delta)$（仍由\cref{thm:exp-convex-ball}）。我们利用这条引理来定义 $Y_{n}$。
\item 若 $\rho_{n}-A<s<\rho_{n}+A$，那么我们利用这一事实：$\ell_{n}(s+\rho_{n})$ 对 $s\in[-A,A]$ \trUniform{}地趋于 $v(s)$。
\item 若 $s>\rho_{n}+A>\tau'_{n}$，那么我们有 $\ell_{n}(s)\in B(z,\delta)$ 与 $w_{\rho_{n}}(s)=v(s-\rho_{n})\in B(z,\delta)$，这就使得我们能够像前面那样利用\cref{thm:exp-convex-ball}来定义 $Y_{n}$。
\end{enumerate}

至此，当 $u$ 与 $v$ 不是关于 $s$ 常值的时，证明结束。若两个解之一（或两者）是常值的，那么证明会更简单。比如说，假定 $v$ 是关于 $s$ 常值的（特别地，此时我们有 $y=z$）。构造 $\ell_{n}$ 的\trBrokenTrajectory{}极限的\cref{thm:broken-limit-decomposition}的证明随后便蕴含：对每个 $s>s^{\star}-\rho_{n}$（某个常数 $s^{\star}$），都有 $\ell_{n}(s)\in B(z,\delta)$。而由于对 $s\ge-\rho_{n}+A$，
\[
u(s+\rho_{n})\in B(y,\delta)=B(z,\delta)
\qquad\text{且}\qquad
v(s-\rho_{n})=y=z\in B(z,\delta)\text{，}
\]
\trPregluing{} $w_{\rho_{n}}(s)$ 显然具有同样的性质。因此，只需对 (1) 与 (2) 中所述的区间定义 $Y_{n}$ 就够了，而这可以用类似的方式来完成。
\end{proof}

\cref{thm:yn-w1p-vanish}的证明与\cref{thm:yn-w1p-null}的证明、以及与其相类似的\cref{thm:yn-w1p}的证明（在\cref{subsec:step3-proof-964}中）完全相同。事实
\[
Y_{n}\in W^{1,p}\bigl(w_{\rho_{n}}^{\star}TW\bigr)
\]
由（\cref{thm:yn-exp-decay} 的）类似物推出（以同样的方式证明）\footnote{译注：原书本式印作 $Y_{n}\in W^{1,p}(w_{\rho}^{\star}TW)$（下标为单 $\rho$）；此处 $Y_{n}$ 与\cref{thm:elln-exp-yn}陈述中的 $w_{\rho_{n}}^{\star}TW$ 相对应，下标疑脱一 $n$，系排印之误，此处径改。——译者注}。

为了证明 $\lim\bigl\|Y_{n}\bigr\|_{W^{1,p}}=0$，我们如同\cref{thm:yn-w1p-null}的证明中那样做一个估计，此时要利用（\cref{thm:taylor-remainder-estimate}的）类似物。注意，\cref{thm:chi-n-w1p}在这里是有用的（因为 $\mu(x)=\mu(y)=\mu(z)$，从而 $W_{n}=0$）；它使得 $Y_{n}$ 的 $W^{1,p}$ 范数的估计更容易得到。\noqed

\begin{proof}[\cref{thm:elln-chi-rho}的证明的结尾]
我们要证明：对充分大的 $n$，
\[
\ell_{n}=\chi_{\rho_{n}}(u,v)\text{，}
\]
换句话说，
\[
\exp_{w_{\rho_{n}}}Y_{n}=\exp_{w_{\rho_{n}}}\gamma_{\rho_{n}}\text{，}
\]
也就是说，$Y_{n}=\gamma_{\rho_{n}}$。我们知道
\[
\scr{F}_{\rho_{n}}\bigl(Y_{n}\bigr)=\scr{F}_{\rho_{n}}\bigl(\gamma_{\rho_{n}}\bigr)=0
\]
并且 $L_{\rho_{n}}=(d\scr{F}_{\rho_{n}})_{0}$ 是\trBijective{}的（正如我们在\cref{subsec:definition-chi}中所见）。但是我们不能直接应用\trLocalInversionTheorem，因为 $\scr{F}_{\rho_{n}}$ 在其上为\trBijective{}的那个邻域的大小可能随 $n$ 增大而减小，而 $Y_{n}$ 未必含于这个邻域之中。我们利用唯一性\cref{thm:general-newton-picard}，它陈述道：借助\trNewton--\trPicard{}方法得到的 $\scr{F}_{\rho_{n}}(Y)=0$ 的解 $\gamma_{\rho_{n}}$ 是球 $B(0,\varepsilon_{0})$ 中唯一的一个，其中常数 $\varepsilon_{0}$ 不依赖于 $n$。

从而，由于 $Y_{n}$ 的 $W^{1,p}$ 范数趋于 $0$，我们有 $Y_{n}=\gamma_{\rho_{n}}$，这就结束了\cref{thm:elln-chi-rho}的证明。
\end{proof}

于是我们便得到了 $\chi_{\rho}$ 对充分大的 $\rho$ 的\trSurjective{}性，这正是结束\cref{thm:chi-bijection}的证明所需要的。\noqed

这就平凡地蕴含了\cref{thm:exists-rho-coincide}：\cref{thm:gamma-rho-regular-large}断言存在任意大的 $\rho$，使得 $\Gamma_{\rho}$ 是\trRegular{}的，而\cref{thm:chi-bijection}蕴含：从 $\rho$ 的某个值起，
\[
\Phi^{\Gamma_{\rho}}=\Phi^{\Gamma''}\circ\Phi^{\Gamma'}
\]
只要左边是有定义的。\footnote{译注：原书此式印作 $\Phi^{\Gamma_{\rho}}=\Phi^{\Gamma'}\circ\Phi^{\Gamma''}$；按两个复合因子各自的定义域与陪域，右端应作 $\Phi^{\Gamma''}\circ\Phi^{\Gamma'}$（与\cref{thm:exists-rho-coincide}陈述处的径改相一致），疑系排印之误，此处径改。——译者注}

最后，\cref{thm:exists-rho-coincide}显然蕴含\cref{thm:composition-regular-homotopy}，后者正是本节的对象。\cref{thm:homology-independent-gamma}与\cref{thm:composition-regular-homotopy}蕴含了\trFloerHomology{}相对于\trRegular{}偶对 $(H,J)$ 选取的\trInvariance。
